\section{Discussion}

\para{Interpretation of Results} Our findings demonstrate that Inductive Moment Matching can be successfully adapted to the reinforcement learning setting. The stability of \immq across different dataset sizes suggests that the particle-based representation, combined with \mmd and inductive bootstrapping, provides a powerful and flexible way to capture return distributions. The observation that a higher inductive penalty $\alpha$ improves \mmd matching while increasing variance is particularly significant. It implies that by explicitly anchoring the first moment (the expected return), we alleviate the optimization burden on the \mmd term, allowing it to more accurately shape the higher-order moments of the distribution. This leads to a more expressive and diverse set of particles, which is vital for downstream tasks like uncertainty-aware exploration.

\para{Comparison to Constrained Methods} Unlike \cfiftyone or \qrdqn, \immq does not require the user to specify categorical supports or fixed quantiles. This unconstrained nature is a major advantage in continuous control, where the range and shape of returns can vary drastically across states and actions. While methods like \iqn provide an implicit quantile representation, \immq offers an explicit particle-based view that may be more intuitive for certain types of risk analysis or exploration strategies (e.g., Thompson sampling over particles).

\para{Limitations} Despite its advantages, \immq has several limitations. The computation of the \mmd loss is $O(N^2)$, where $N$ is the number of particles. While $N=32$ was computationally manageable in our experiments, larger particle sets required for highly complex distributions could become a bottleneck. Furthermore, our evaluation was limited to an offline setting using offline policy evaluation (OPE). The true test of \immq's sample efficiency lies in its integration into an online actor-critic loop, where the learned distribution can actively guide exploration and policy improvement.
